%

\subsection{Stożki} % lang-pl
\subsection{Coni} % lang-it
Stożki. % lang-pl
Coni. % lang-it

Jerzy Konarski opisze w $\Delta_{93}^{5}$ powierzchnie rozwijalne, czyli takie, które można otrzymać z kartki papieru: poza płaszczyzną są to stożek, walec i powierzchnia stycznych do krzywej przestrzennej; wśród powierzchni prostokreślnych rozpozna je po tym, że wektory normalne wzdłuż tworzącej są równoległe, albo po zerowej krzywiźnie Gaussa. % lang-pl
Jerzy Konarski descriverà in $\Delta_{93}^{5}$ le superfici sviluppabili, cioè quelle che si possono ottenere da un foglio di carta: oltre al piano sono il cono, il cilindro e la superficie delle tangenti a una curva spaziale; tra le superfici rigate le riconoscerà dal fatto che i vettori normali lungo una generatrice sono paralleli, oppure dalla curvatura gaussiana nulla. % lang-it

%